\subsection{When the chain does not share one algorithm}
\label{sec:mixedchain}

Every chain so far uses one algorithm from leaf to root, and no migration
looks like that. A certificate carries two algorithms: the key it holds and
the signature its issuer put on it. Those come apart during a transition, and
in a particular order, because leaf certificates are reissued in days,
intermediates live for years, and roots are pinned in trust stores and move
last.

\begin{table}[t]
\caption{Chains whose levels migrate at different times, leaf first and root last; the root signs the intermediate but is not carried. \emph{Signer PQ} is whether every load-bearing signature in the signer's path is post-quantum. The last two columns add an RFC~3161 timestamp: with a classical TSA the time evidence stays forgeable, so a quantum adversary can backdate the manifest even when the signer's path is post-quantum.}
\label{tab:mixedchain}
\small
\setlength{\tabcolsep}{3.5pt}
\begin{tabular}{llrcrr}
\toprule
 & Chain & COSE & Signer & \multicolumn{2}{c}{+ timestamp (B)}\\
\cmidrule(lr){5-6}
Stage & leaf/int./root & (B) & PQ & ES256 TSA & ML-DSA-65 TSA\\
\midrule
Classical & ES256/ES256/ES256 & 1,585 & no & 3,651 & 17,209\\
Leaf migrated & D65/ES256/ES256 & 6,724 & no & 8,788 & 22,348\\
Leaf and intermediate & D65/D65/ES256 & 11,867 & no & 13,933 & 27,491\\
PKI first, leaves last & ES256/D65/D65 & 9,971 & no & -- & --\\
Fully migrated & D65/D65/D65 & 15,113 & yes & 17,179 & 30,737\\
Assurance ladder & D44/D65/D87 & 14,902 & yes & -- & --\\
Ladder inverted & D87/D65/D44 & 16,182 & yes & -- & --\\
Cheap leaves, strong root & D44/D44/D87 & 13,373 & yes & -- & --\\
\bottomrule
\end{tabular}
\end{table}


Table~\ref{tab:mixedchain} separates two things usually spoken of as one.
Bytes are set by which levels have migrated. Quantum resistance is set by the
weakest load-bearing signature: the manifest signature, the signature on the
leaf certificate and the signature on the intermediate. The root's
self-signature is not load-bearing, and the root is not carried, but its
algorithm still matters because it signs the intermediate. A post-quantum
leaf under the existing classical PKI costs \NmixLeafMig{} times the
classical structure, \NmixPctFullLeafMig\% of what full migration costs;
migrating the intermediate too brings it to \NmixLeafInt{} times. Neither
makes the signer's path quantum resistant, because an adversary who can forge
the issuer's ECDSA signature mints a leaf certificate and signs whatever it
likes.

Migrating the PKI first while issuing classical leaves for compatibility is
the order an operator reaches by reasoning about disruption rather than about
adversaries. It costs \NmixPkiFirst{} times the classical structure,
\NmixPctFullPkiFirst\% of the full price, and leaves the manifest itself
forgeable, since its signature is still ECDSA. What it does buy is a
post-quantum trust infrastructure ready to issue post-quantum leaves without a
second change at the anchor, so its value lies in what follows it, not in any
protection it gives the manifests signed meanwhile.

The timestamp needs its own migration. With a classical time-stamping
authority, a fully migrated signer's manifest can still be backdated by an
adversary who forges the timestamp, which undoes the argument a timestamp
exists to make, that the manifest was signed before its certificate was
revoked. A post-quantum timestamp costs \NmixWholeFullPq{} bytes of manifest
cryptography against \NmixWholeFullEs{} with an ES256 authority, and the
provenance is quantum resistant only when both are post-quantum.

An assurance ladder, a category-2 leaf under a category-3 intermediate under a
category-5 root, puts the higher margin on the longer-lived key and costs
\NmixLadder{} times the classical structure, against \NmixFull{} for a uniform
category-3 chain: the strongest key sits in the root, which is not carried.
Inverting the ladder costs \NmixInverted{} times and is the configuration to
check for in review, since it puts the weakest parameter set on the key that
lives longest and is hardest to replace.
